在圆维签名的口径下，有限 torsion 账本既可作为实现层的外置寄存器，也可由内禀算术数据以规范方式生成。
以下给出一条以判别群（cokernel）为核心的构造链：对称整双线性型的判别群与其对偶紧环面同态的覆盖核自然同构；对数域迹格，该判别群等同于 codifferent 商并由判别式决定坏素支撑；局部 Smith 正规形进一步给出按素数分解的“轴数—深度”离散谱；最后，判别账本携带的 \(\QQ/\ZZ\)-值配对诱导规范有限傅里叶算子，并与素分解严格张量化兼容。

\begin{theorem}[判别群与覆盖核：最小有限账本的规范实现]\label{thm:xi-disc-ledger-kernel-dual}
\leanverified{paper\_xi\_disc\_ledger\_kernel\_dual}
设 \(M\) 为秩 \(n\) 的自由 \(\ZZ\)-模，\(B:M\times M\to \ZZ\) 为对称双线性型，且其 \(\QQ\)-线性延拓在 \(M_{\QQ}:=M\otimes_{\ZZ}\QQ\) 上非退化。
定义同态
\[
\varphi_B:M\to M^\vee:=\operatorname{Hom}_{\ZZ}(M,\ZZ),\qquad
x\mapsto (y\mapsto B(x,y)),
\]
并定义判别群
\[
D_B:=\operatorname{coker}(\varphi_B)=M^\vee/\varphi_B(M).
\]
令紧环面
\[
\TT_M:=\operatorname{Hom}(M,\RR/\ZZ)\cong (\RR/\ZZ)^n.
\]
则：
\begin{enumerate}
  \item \(\varphi_B\) 在 \(\QQ\)-上线性同构，故 \(\widehat{\varphi_B}:\TT_M\to \TT_M\) 为有限覆盖（isogeny），并存在自然同构
  \[
  \ker(\widehat{\varphi_B})\ \cong\ \widehat{D_B}:=\operatorname{Hom}(D_B,\RR/\ZZ),
  \]
  从而 \(|\ker(\widehat{\varphi_B})|=|D_B|\)。
  \item 设 \(R\) 为有限生成离散交换群且 \(r:\TT_M\to \widehat R\) 为连续群同态，使得
  \[
  (\widehat{\varphi_B},r):\TT_M\longrightarrow \TT_M\times \widehat R
  \]
  为单射，则 \(\widehat R\) 必含一个与 \(\ker(\widehat{\varphi_B})\) 同构的有限子群；等价地，\(R\) 的 torsion 商中必出现与 \(D_B\) 同阶的因子。特别地，任何此类实现的不可规避素数支撑包含 \(|D_B|\) 的全部素因子（等价包含 \(|\det(B)|\) 的全部素因子）。
\end{enumerate}
\end{theorem}
\begin{proof}
取 \(\ZZ\)-基使 \(M\cong \ZZ^n\)。此时 \(\varphi_B\) 由整矩阵 \(A\in M_n(\ZZ)\) 表示，且
\(
D_B\cong \ZZ^n/A\ZZ^n
\)。
对短正合列
\[
0\to \ZZ^n \xrightarrow{A} \ZZ^n \to \ZZ^n/A\ZZ^n\to 0
\]
取 Pontryagin 对偶（离散群的对偶函子保持正合），得
\[
0\to \widehat{\ZZ^n/A\ZZ^n}\to \widehat{\ZZ^n}\xrightarrow{\widehat A}\widehat{\ZZ^n}\to 0.
\]
又 \(\widehat{\ZZ^n}\cong (\RR/\ZZ)^n\cong \TT_M\) 且 \(\widehat A\) 即由 \(A^{\mathsf T}\) 在 \((\RR/\ZZ)^n\) 上诱导的覆盖同态；由于 \(B\) 对称，可在更换基后取 \(A=A^{\mathsf T}\)，从而得到 (1)。
\par
对 (2)，若 \((\widehat{\varphi_B},r)\) 单射，则 \(r\) 在 \(\ker(\widehat{\varphi_B})\) 上亦为单射，故 \(\ker(\widehat{\varphi_B})\) 可嵌入 \(\widehat R\)。
而任意紧交换群 \(\widehat R\) 的有限子群对偶于 \(R\) 的 torsion 商，因此 \(R\) 的 torsion 部分必须包含与 \(\ker(\widehat{\varphi_B})\) 同构的对偶信息；由 (1) 得其阶为 \(|D_B|\)，素数支撑断言随之成立。
\end{proof}

\begin{theorem}[迹格的判别账本与坏素支撑]\label{thm:xi-trace-lattice-discriminant-ledger}
令 \(K/\QQ\) 为次数 \(n\) 的数域，\(\mathcal O_K\) 为其整数环。
取迹配对
\[
B(x,y):=\operatorname{Tr}_{K/\QQ}(xy),\qquad x,y\in\mathcal O_K,
\]
并定义迹对偶格（codifferent）
\[
\mathcal O_K^\vee:=\{\alpha\in K:\operatorname{Tr}_{K/\QQ}(\alpha\,\mathcal O_K)\subseteq \ZZ\}.
\]
令
\[
D_K:=\mathcal O_K^\vee/\mathcal O_K.
\]
则存在自然同构 \(D_K\cong D_B\)，并且
\[
|D_K|=|\Delta_K|,
\]
其中 \(\Delta_K\) 为 \(K\) 的绝对判别式。
从而 \(\widehat{\varphi_B}:\TT_{\mathcal O_K}\to\TT_{\mathcal O_K}\) 的覆盖核同构于 \(\widehat{D_K}\)，且任何使该同态在扩展保存口径下可逆的审计实现都不可避免地引入坏素支撑
\[
\{p:\ p\mid \Delta_K\}.
\]
\end{theorem}
\begin{proof}
对任意整基 \((e_i)\) 的迹矩阵 \(P=(\operatorname{Tr}(e_ie_j))\) 有 \(\det(P)=\Delta_K\)。
而 \([\mathcal O_K^\vee:\mathcal O_K]=|\det(P)|\) 为经典格论恒等式，故 \(|D_K|=[\mathcal O_K^\vee:\mathcal O_K]=|\Delta_K|\)。
同构 \(D_K\cong D_B\) 来自 \(\varphi_B\) 与对偶格的定义对齐。
最后的坏素支撑结论由定理 \ref{thm:xi-disc-ledger-kernel-dual}(2) 立即推出。
\end{proof}

\begin{proposition}[局部最小 \(p\)-进轴数与迹矩阵秩亏]\label{prop:xi-disc-ledger-local-axis-rank-defect}
沿用定理 \ref{thm:xi-trace-lattice-discriminant-ledger} 的记号，取 \(\mathcal O_K\) 的任意整基并令 \(P=(\operatorname{Tr}(e_ie_j))\in M_n(\ZZ)\)。
对每个素数 \(p\) 定义
\[
g_p(K):=\dim_{\FF_p}\bigl(D_K/pD_K\bigr).
\]
则有显式公式
\[
g_p(K)=n-\operatorname{rank}_{\FF_p}(P\bmod p).
\]
并且 \(g_p(K)\) 等于 \(D_K\) 的 \(p\)-主分量所需的最小生成元个数，从而给出任何审计实现中不可规避的独立 \(p\)-进外置轴数量下界。
\end{proposition}
\begin{proof}
由 \(D_K\cong \ZZ^n/P\ZZ^n\) 得
\[
D_K/pD_K\cong (\ZZ^n/P\ZZ^n)\otimes_{\ZZ}\FF_p\cong \FF_p^n/(P\bmod p)\FF_p^n,
\]
故其维数为 \(n-\operatorname{rank}_{\FF_p}(P\bmod p)\)。
有限阿贝尔群 \(G\) 的 \(p\)-秩 \(r_p(G):=\dim_{\FF_p}(G/pG)\) 等于其 \(p\)-主分量的最小生成元个数，代入 \(G=D_K\) 即得结论。
\end{proof}

\begin{theorem}[判别素账本的轴数—深度全分解]\label{thm:xi-disc-ledger-axis-depth-smith}
\leanverified{paper\_xi\_disc\_ledger\_axis\_depth\_smith}
沿用上述记号。
对每个素数 \(p\)，令 \(P\) 在 \(\ZZ_p\) 上的 Smith 正规形为
\[
UPV=\operatorname{diag}(p^{\alpha_{p,1}},\dots,p^{\alpha_{p,n}}),
\qquad
U,V\in \operatorname{GL}_n(\ZZ_p),
\qquad
0\le \alpha_{p,n}\le\cdots\le \alpha_{p,1},
\]
并令
\[
g_p(K):=\#\{i:\alpha_{p,i}>0\},\qquad
\boldsymbol{\alpha}_p(K):=(\alpha_{p,1}\ge\cdots\ge\alpha_{p,g_p(K)}\ge 1).
\]
则：
\begin{enumerate}
  \item \((D_K)_{(p)}\cong \bigoplus_{i=1}^{g_p(K)}\ZZ/p^{\alpha_{p,i}}\ZZ\)，并且序列 \(\boldsymbol{\alpha}_p(K)\) 与整基选择无关；
  \item 判别式的 \(p\)-进指数满足恒等式
  \[
  v_p(\Delta_K)=\sum_{i=1}^{g_p(K)}\alpha_{p,i}.
  \]
\end{enumerate}
因此 \(g_p(K)\) 给出不可规避的 \(p\)-进外置轴数目，而 \(\boldsymbol{\alpha}_p(K)\) 给出每一轴的必需深度；二者共同构成判别素账本的完备离散谱。
\end{theorem}
\begin{proof}
Smith 正规形在 PID \(\ZZ_p\) 上存在且不变量唯一，给出 cokernel 的直和分解与基不变性。
另一方面，对角元之积等于 \(\det(P)\)，而 \(\det(P)=\Delta_K\)，故
\(
v_p(\Delta_K)=\sum_{i=1}^{n}\alpha_{p,i}
=\sum_{i=1}^{g_p(K)}\alpha_{p,i}
\)。
\end{proof}

\begin{corollary}[二次数域的循环判别账本]\label{cor:xi-disc-ledger-quadratic-cyclic}
若 \([K:\QQ]=2\)，则对每个坏素 \(p\mid \Delta_K\) 有 \(g_p(K)=1\) 且
\[
(D_K)_{(p)}\cong \ZZ/p^{v_p(\Delta_K)}\ZZ.
\]
从而 \(D_K\cong \ZZ/|\Delta_K|\ZZ\) 为循环群；审计意义下每个坏素仅需单一 \(p\)-进外置轴，其深度恰为 \(v_p(\Delta_K)\)。
\end{corollary}
\begin{proof}
二次情形下 \(P\in M_2(\ZZ)\) 且 \(\det(P)\neq 0\)。其 Smith 不变量在同构意义下为 \((1,|\det(P)|)\)，故 cokernel 循环；局部化到 \(\ZZ_p\) 后仅出现一个非零指数 \(v_p(\det(P))=v_p(\Delta_K)\)。
\end{proof}

\begin{proposition}[判别账本的规范有限傅里叶算子与素分解张量化]\label{prop:xi-disc-ledger-fourier-tensorization}
令 \(D\) 为有限阿贝尔群，\(\langle\cdot,\cdot\rangle:D\times D\to \QQ/\ZZ\) 为非退化双线性配对，并定义双角色
\[
\chi(a,b):=\exp\!\bigl(2\pi \mathrm{i}\,\langle a,b\rangle\bigr).
\]
在群代数 \(\CC[D]\) 上定义算子
\[
(\mathcal F_\chi f)(a):=|D|^{-1/2}\sum_{b\in D}\chi(a,b)\,f(b).
\]
则：
\begin{enumerate}
  \item \(\mathcal F_\chi\) 为酉算子，且仅由配对 \((D,\langle\cdot,\cdot\rangle)\) 决定；
  \item 若 \(D=\bigoplus_{p}D_{(p)}\) 为 \(p\)-主分解且配对按该直和正交分解，则存在规范张量同构
  \[
  \CC[D]\cong \bigotimes_p \CC[D_{(p)}],
  \qquad
  \mathcal F_\chi\cong \bigotimes_p \mathcal F_{\chi_{(p)}},
  \]
  其中 \(\chi_{(p)}\) 为 \(\chi\) 在 \(D_{(p)}\) 上的限制。
\end{enumerate}
在定理 \ref{thm:xi-trace-lattice-discriminant-ledger} 的设定下，判别账本 \(D_K\) 上存在由迹对偶格诱导的自然 \(\QQ/\ZZ\)-值配对，从而 \(\mathcal F_{\chi_K}\) 的谱结构严格张量化分解为各坏素的局部谱因子。
\end{proposition}
\begin{proof}
对任意 \(a,a'\in D\)，利用非退化性可得字符正交关系
\[
\sum_{b\in D}\chi(a,b)\overline{\chi(a',b)}=
\sum_{b\in D}\chi(a-a',b)=
\begin{cases}
|D|,& a=a',\\
0,& a\neq a',
\end{cases}
\]
从而矩阵 \((|D|^{-1/2}\chi(a,b))\) 为酉矩阵，即得 (1)。
对 (2)，当 \(D=\bigoplus_p D_{(p)}\) 且配对正交分解时，双角色按素因子化：\(\chi=\prod_p \chi_{(p)}\)。
将求和按直积拆分并用 \(|D|^{1/2}=\prod_p |D_{(p)}|^{1/2}\) 得到算子的张量分解，群代数的张量同构由有限直和的标准基给出。
最后一段为定义对齐与前述分解的直接应用。
\end{proof}

\begin{conjecture}[判别账本相位谱的四点均匀性]\label{conj:xi-disc-ledger-fourier-spectrum-fourpoint}
在命题 \ref{prop:xi-disc-ledger-fourier-tensorization} 的设定下，进一步假设配对 \(\langle\cdot,\cdot\rangle\) 对称。
则 \(\mathcal F_\chi^4=\mathrm{Id}\)，从而 \(\mathcal F_\chi\) 的谱包含于 \(\{\pm 1,\pm \mathrm{i}\}\)。
固定次数 \(n\ge 2\)，令 \(K\) 在所有 \([K:\QQ]=n\) 且 \(|\Delta_K|\to\infty\) 的数域族中变化，并令 \(\mu_K\) 为 \(\mathcal F_{\chi_K}\) 的经验谱测度（按特征值计重）。
猜想 \(\mu_K\) 在弱-\(\ast\) 意义下趋于 \(\{\pm 1,\pm \mathrm{i}\}\) 上的均匀测度。
\end{conjecture}
